\section{Scope and notation}
\label{sec:notation}

These notes are Part~0 of the GeoFDI theory series. They fix notation and the
concrete form of the symmetry representations, state the standing assumptions
\ref{as:dyn}--\ref{as:erg} together with $\varepsilon$-approximate variants,
audit procedures, and degradation paths, and formalize the null hypothesis
$H_0$ (distributional symmetry of healthy gait data), the structured
alternative $H_1$, and the recalibrated null $H_0'$. Test construction,
cross-cycle exchangeability, the isotypic detectability characterization, and
full proofs are deferred to later parts (workstreams N1-1 through N1-4).

The project sits between two literatures. Model-based proprioceptive fault and
collision detection builds residuals from a dynamics model and monitors them
\cite{haddadin2017robot,evangelisti2024datadriven}; symmetry-aware robotics
identifies the finite morphological symmetry group of a robot and exploits it
as an algebraic structure on models and data
\cite{ordonez2023discrete,ordonez2025morphological}. GeoFDI uses the second
structure to stand in for the first's model: under the assumptions of
Section~\ref{sec:assumptions}, \emph{healthy} closed-loop gait data is
distribution-invariant under the gait's spatio-temporal symmetry group, so
fault detection reduces to statistical testing of group invariance --- a
hypothesis class that admits exact, distribution-free permutation tests
\cite{hemerik2018exact} and anytime-valid sequential e-processes
\cite{perezortiz2024estatistics}.

\subsection{State, input, output, and residual spaces}
\label{sec:spaces}

We consider a floating-base legged robot with $n_j$ actuated joints. The
configuration space is
\[
  Q \;=\; \SE(3) \times \mathbb{T}^{n_j},
  \qquad q = (q_B, q_J),
\]
with $q_B = (R_{B}, p_{B}) \in \SE(3)$ the base pose (body-to-world rotation
and world position) and $q_J \in \mathbb{T}^{n_j}$ the joint coordinates. The
state is $x = (q, v) \in X = TQ$, where $v = (\xi, \dot q_J) \in \R^{6+n_j}$
stacks the body twist $\xi = (v_b, \omega_b)$ and the joint rates. Actuation
commands (desired torques or motor commands) live in $u \in U \subseteq
\R^{n_u}$. The measured output $y \in Y$ collects the proprioceptive channels:
IMU specific force $a$ and angular rate $\omega$, joint encoders and rates,
motor currents or measured torques, and per-leg contact indicators
$c \in \{0,1\}^{4}$.

A \emph{residual} $r \in R$ is any causally computed function of the data
stream designed to be small or structured when the robot is healthy. The
theory below treats $y$ and $r$ generically; the concrete channels planned for
GeoFDI are the innovation sequences of a contact-aided invariant extended
Kalman filter \cite{hartley2020contact,barrau2017invariant} and
generalized-momentum residuals \cite{evangelisti2024datadriven}, alongside raw
tracking errors. When a residual generator is built equivariantly, the induced
representation $\repR$ on $R$ makes the residual stream inherit the symmetry
of the data (workstream N2).

\subsection{Morphological symmetry group and representations}
\label{sec:group}

$G$ denotes the finite morphological symmetry group of the robot: the group of
transformations that map the mechanism onto itself (relabeling mirrored limbs
and reflecting spatial quantities) while preserving its kinematic and dynamic
description \cite{ordonez2023discrete,ordonez2025morphological}. $G$ acts on
every space in play through representations
\[
  \repQ \text{ on } Q,\qquad
  \repX \text{ on } X,\qquad
  \repU \text{ on } U,\qquad
  \repY \text{ on } Y,\qquad
  \repR \text{ on } R,
\]
whose concrete matrices are spelled out in Section~\ref{sec:reps}. A map
$h : X \times U \to X$ is \emph{$G$-equivariant} if
$h(\repX(g)x, \repU(g)u) = \repX(g)\,h(x,u)$ for all $g \in G$. We write
$W = X \times U \times Y$ for the per-sample data space and
$\repW = \repX \oplus \repU \oplus \repY$ for the block action on it.
Exogenous environment inputs (terrain heights, pushes) are denoted
$w$ with action $\repenv$.

\begin{table}[ht]
\centering
\small
\begin{tabular}{@{}ll@{}}
\toprule
symbol & meaning \\
\midrule
$Q = \SE(3)\times\mathbb{T}^{n_j}$ & configuration space: floating base $\times$ joints \\
$x = (q,v) \in X$ & state; $v = (\xi, \dot q_J)$ body twist and joint rates \\
$u \in U$ & actuation commands (torques / motor commands) \\
$y \in Y$ & proprioceptive outputs (IMU, encoders, currents, contacts) \\
$r \in R$ & residual channels (InEKF innovations, momentum residuals, \dots) \\
$G$, $\gs$ & finite morphological symmetry group; sagittal reflection ($G=\mathbb{C}_2$ for M1) \\
$\repQ,\repX,\repU,\repY,\repR$ & representations of $G$ on the respective spaces \\
$W$, $\repW$ & data space $X\times U\times Y$ and its block action \\
$w$, $\repenv$ & exogenous environment inputs and their action \\
$\theta \in S^1 = \R/\mathbb{Z}$ & gait phase (period normalized to 1); $T$ = period [s] \\
$\symgrp \subset G \times S^1$ & spatio-temporal symmetry group of the commanded gait \\
$\sigma = (g, \Delta\theta)$ & element of $\symgrp$: spatial action paired with phase shift \\
$Z_k$ & phase-registered per-cycle data element, cycle $k$ \\
$\rho(\sigma)$ & action on cycle elements: $(\rho(\sigma)Z)(\theta) = \repW(g)\,Z(\theta+\Delta\theta)$ \\
$\eps{dyn},\eps{ctrl},\eps{env},\eps{noise}$ & per-transition equivariance defects (A1$'$--A4$'$) \\
$\epsbar{dyn},\epsbar{ctrl},\epsbar{env},\epsbar{noise}$ & cycle-level defects (Definition~\ref{def:cycle-eps}) \\
$\Tc$ & closed-loop transition steps per gait cycle \\
$\nu$, $\nu_0$ & asymmetry functional and its calibration value (for $H_0'$) \\
\bottomrule
\end{tabular}
\caption{Symbols used throughout the GeoFDI notes.}
\label{tab:symbols}
\end{table}

\subsection{Gait phase and per-cycle data elements}
\label{sec:cycles}

The commanded gait is periodic with period $T$; the \emph{gait phase}
$\theta(t) \in S^1$ is the normalized position within the cycle, estimated from
contact events or a phase oscillator (workstream N1/N3; the phase-registration
error budget enters the $\varepsilon$ accounting of
Section~\ref{sec:h0prime}). Slicing the telemetry at phase zero yields the
\emph{per-cycle data element}
\[
  Z_k \;=\; \bigl\{\, (x(t), u(t), y(t)) \;:\; t \in \text{cycle } k \,\bigr\},
\]
which after phase registration is a random function
$Z_k : S^1 \to W$ --- the fundamental statistical object of this project: all
tests act on the sequence $(Z_k)_{k\ge 1}$, never on raw time series.%
\footnote{In implementation, channels are resampled on a fixed $N$-point phase
grid, so $Z_k$ is a random vector in $\R^{dN}$ ($d$ channels); every law below
lives on that finite-dimensional space, which sidesteps measure-theoretic
subtleties on function spaces. $N$ is a protocol parameter.}

For $\sigma = (g, \Delta\theta) \in G \times S^1$ define the action on cycle
elements by
\begin{equation}
  \label{eq:cycle-action}
  (\rho(\sigma) Z)(\theta) \;=\; \repW(g)\, Z(\theta + \Delta\theta),
\end{equation}
i.e.\ spatially transform every channel and shift the phase. This is a group
action, and restricting it to the gait's symmetry group $\symgrp$
(Section~\ref{sec:gait}) gives the invariance tested by $H_0$.

\subsection{Concretization for the M1 quadruped}
\label{sec:m1}

For the M1 quadruped we take
\[
  G \;=\; \mathbb{C}_2 \;=\; \{e, \gs\},
\]
where $\gs$ is the sagittal (left--right) reflection, acting on leg labels by
the permutation
\[
  \sigma_{\mathrm{s}} \;=\; (\mathrm{LF}\;\mathrm{RF})(\mathrm{LH}\;\mathrm{RH}),
\]
i.e.\ left-front $\leftrightarrow$ right-front and left-hind
$\leftrightarrow$ right-hind. Front--hind symmetry is \emph{not} assumed (mass
distribution and leg geometry differ between front and hind pairs), so the
richer groups $\mathbb{C}_2\times\mathbb{C}_2$ or $\mathbb{K}_4$ of idealized
quadrupeds are unavailable.

\paragraph{Joint count, ordering and reflection signs (resolved on hardware).}
The M1 platform is a \emph{wheeled-legged} quadruped: $n_j = 16$, four legs
$\times$ (ABAD = hip roll, HIP = hip pitch, KNEE = knee pitch, WHEEL). The
ordering and every reflection sign in Table~\ref{tab:joint-signs} were fixed on
the 2026-08-10 hardware corpus (\texttt{docs/protocol/m1\_h\_data\_audit.md};
telemetry names \texttt{fl1\_hip\_roll \dots br4\_foot}), superseding the
placeholder that stood here.

Two conventions must be kept apart. In the \emph{uniform-axis} convention used
throughout this document, a channel's mirror sign is decided by whether it is a
polar or an axial quantity about the sagittal plane: the roll-axis joint flips,
the pitch-axis joints do not. In the \emph{vendor} frame the manufacturer may
already define mirrored positive directions, in which case the raw telemetry
carries extra sign flips that the loader must undo before the representation
$\repW(\gs)$ of \eqref{eq:cycle-action} applies --- this is the case for M1,
where all four joints flip between mirror legs in vendor coordinates.

\begin{table}[ht]
\centering
\footnotesize
\begin{tabular}{@{}p{2.6cm}p{3.5cm}p{2.5cm}cc@{}}
\toprule
platform & channel & kind & $\repW(\gs)$ & vendor$\to$uniform (R legs) \\
\midrule
M1 (16 joints) & ABAD (hip roll) & axial-about-$x$ & $-1$ & $+1$ \\
               & HIP (hip pitch) & pitch axis & $+1$ & $-1$ \\
               & KNEE (knee pitch) & pitch axis & $+1$ & $-1$ \\
               & WHEEL (rate) & pitch axis & $+1$ & $-1$ \\
\midrule
Go2 (no joints) & \texttt{foot\_pos}/\texttt{foot\_vel} $(x,y,z)$ & polar vector & $(+1,-1,+1)$ & --- \\
               & \texttt{foot\_force} & magnitude & $+1$ & --- \\
\midrule
both & IMU accelerometer & polar & $\operatorname{diag}(+1,-1,+1)$ & --- \\
     & IMU gyroscope & axial & $\operatorname{diag}(-1,+1,-1)$ & --- \\
\bottomrule
\end{tabular}
\caption{Per-channel reflection signs under $\gs$, verified on hardware. M1:
Sprint~8 audit; the wheel \emph{angle} is unbounded (the vendor wraps it to
$[-\pi,\pi)$) and is therefore excluded from the data element, while the wheel
\emph{rate} and torque stay. The rightmost column is the extra flip the loader
applies to the right legs to leave the vendor frame; the forward-rolling sign of
the wheel is additionally $-1$ on the left legs and $+1$ on the right. Go2
(Sprint~9): the high-level API exposes \emph{no} joint stream, so the element is
built from the robot's own forward kinematics
(\texttt{foot\_position\_body}/\texttt{foot\_speed\_body}) plus foot force and
the IMU; the same $\mathbb{C}_2$ representation applies channel-wise.}
\label{tab:joint-signs}
\end{table}

For the Go2 corpus of Sprint~9 the situation is different in kind rather than in
degree: the recorded high-level interface (\texttt{SportModeState}) carries no
\texttt{motorState}, so $n_j$ is not observable at all. What it does publish is
the manufacturer's own forward kinematics of each leg, which is a bijective
re-coordinatisation of the leg joint vector, so the mirror pairing and the
isotypic split of Section~\ref{sec:reps} go through unchanged on those channels.
Only the torque channel --- and with it the model-residual ($R^+$) layer --- is
lost.
